\documentclass[UTF8]{ctexart}
\usepackage{geometry}
\geometry{a4paper, margin=2.5cm}
\usepackage{listings}
\usepackage{xcolor}
\usepackage{fontspec}
\usepackage{amsmath}
\usepackage{graphicx}
\usepackage{silence}
\usepackage{lmodern}
\WarningFilter{latex}{Font shape}
\WarningFilter{LaTeX}{Font shape}
\setmonofont{DejaVu Sans Mono}
\lstset{
    basicstyle=\ttfamily\small,
    breaklines=true,
    frame=single,
    numbers=left,
    numberstyle=\tiny,
    tabsize=4,
    language=C,
    keywordstyle=\color{blue},
    commentstyle=\color{green!60!black},
    stringstyle=\color{red}
}

\title{基于哈夫曼编码的文件压缩与解压工具实验报告}
\author{}

\begin{document}
\maketitle

\begin{center}
    \textbf{班级：}2024217802 \quad \textbf{姓名：}严子竣 \quad \textbf{学号：}2024212622
\end{center}

\section{需求分析}
在计算机系统中，无损压缩技术广泛用于减少文件存储空间和传输带宽。哈夫曼编码作为一种经典的基于字符出现频率的最优前缀编码，被应用于 ZIP、PNG、JPEG 等众多文件格式的压缩环节。本实验要求基于哈夫曼编码实现一个通用的文件压缩与解压工具，支持对任意格式的文件进行压缩，生成真正的二进制压缩文件，并能够将压缩文件完全还原为原始文件。

\subsection{程序功能}
\begin{itemize}
    \item \textbf{压缩功能}：
    \begin{enumerate}
        \item 以二进制方式读取源文件，统计所有字节值（0--255）的出现频率。
        \item 根据频率构建哈夫曼树，生成每个字节值的变长编码。
        \item 将原始文件名、原始文件大小、频率表等元数据与编码后的比特流一并写入压缩文件，形成可自解压的二进制存档。
    \end{enumerate}
    \item \textbf{解压功能}：
    \begin{enumerate}
        \item 从压缩文件中提取魔数、文件名、原文件大小、频率表。
        \item 利用频率表重建与压缩时完全一致的哈夫曼树。
        \item 按位读取压缩数据，沿哈夫曼树解码并逐字节写回目标文件，还原出完整原始文件。
    \end{enumerate}
    \item \textbf{错误处理}：对文件打开失败、损坏的压缩文件、内存不足等情况均能输出清晰的错误提示并安全退出。
\end{itemize}

\subsection{输入形式与值的范围}
\begin{itemize}
    \item \textbf{压缩阶段}：
        \begin{itemize}
            \item 通过命令行传入源文件路径和压缩文件路径，或在程序内部交互式提示输入。
            \item 源文件可为任意格式、任意大小（受限于操作系统文件系统和内存）。
            \item 压缩文件名可自定义，推荐以 \texttt{.huff} 为后缀。
        \end{itemize}
    \item \textbf{解压阶段}：
        \begin{itemize}
            \item 通过命令行传入压缩文件路径和输出目录（可选），或交互式输入。
            \item 压缩文件必须符合程序定义的格式，否则提示格式错误。
            \item 解压输出的目录若不存在，程序会报错并退出。
        \end{itemize}
\end{itemize}

\subsection{输出形式}
\begin{itemize}
    \item 压缩成功时输出：\texttt{压缩完成：源文件 -> 压缩文件}
    \item 解压成功时输出：\texttt{解压完成：压缩文件 -> 输出文件}
    \item 发生错误时输出具体原因，例如：\texttt{无法打开输入文件 xxx: Permission denied} 或 \texttt{文件格式错误，不是有效的 Huffman 压缩文件。}
\end{itemize}

\subsection{测试数据}
\begin{itemize}
    \item \textbf{正确输入}：一个普通文本文件、一张小图片或一个二进制文件。
    \item \textbf{边界输入}：空文件（应提示无法压缩）、仅含单一字节的文件（压缩后应能正确还原）。
    \item \textbf{错误输入}：不存在的源文件、无写入权限的输出路径、损坏的压缩文件（缺少魔数、频率表等）。
\end{itemize}

\section{概要设计}

\subsection{问题解决思路}
程序分为压缩和解压两个独立的可执行文件，各自采用以下策略：
\begin{enumerate}
    \item \textbf{压缩流程}：
        \begin{itemize}
            \item 分块读取源文件，用大小为 256 的桶统计各字节出现次数，得到频率数组。
            \item 利用最小堆构建哈夫曼树：为每个出现过的字节创建叶子节点，按频率进行堆排序；每次取出两个最小频率节点合并为新父节点，直到只剩一个根节点。
            \item 递归遍历哈夫曼树，为每个叶子生成二进制编码（向左为 0，向右为 1），编码位数最多 255 位，存储为定长字节数组。
            \item 写入压缩文件头：4 字节魔数 \texttt{"HUFF"}、2 字节原文件名长度、原文件名字符串、8 字节原始文件大小、256×8 字节频率表。
            \item 重新读取源文件，根据编码表将每个字节转换为变长比特流，通过位缓冲逐字节写入压缩文件。
        \end{itemize}
    \item \textbf{解压流程}：
        \begin{itemize}
            \item 读取并验证魔数，确认文件格式。
            \item 读取文件名、原始大小、频率表。
            \item 根据频率表再次利用最小堆重建哈夫曼树，保证与压缩时完全一致。
            \item 按位读取压缩数据，从根节点开始，每读取一个比特决定向左或向右移动，到达叶子节点时输出该字节，然后回到根节点继续，直到还原全部原始大小字节。
        \end{itemize}
\end{enumerate}

\subsection{数据结构定义}
\begin{itemize}
    \item \textbf{哈夫曼树节点}：
\begin{lstlisting}
typedef struct Node {
    int symbol;               // 字节值0~255，内部节点为-1
    uint64_t freq;
    struct Node *left, *right;
} Node;
\end{lstlisting}
    \item \textbf{最小堆}：
\begin{lstlisting}
typedef struct {
    Node **array;
    int size, capacity;
} MinHeap;
\end{lstlisting}
    解压器中，最小堆元素为结构体对 \texttt{(Node*, freq)}，定义形式稍异，但逻辑等价。
    \item \textbf{编码表项}：
\begin{lstlisting}
typedef struct {
    unsigned char code[CODE_BYTES];  // 存储二进制编码，每字节存8位
    int len;                         // 有效位数
} CodeTable;
\end{lstlisting}
    \item \textbf{位写入器（压缩用）}：
\begin{lstlisting}
typedef struct {
    FILE *fp;
    unsigned char buf;
    int pos;                // 当前缓冲位0~7
} BitWriter;
\end{lstlisting}
    \item \textbf{位读取器（解压用）}：
\begin{lstlisting}
typedef struct {
    FILE *fp;
    unsigned char buf;
    int pos;                // 当前剩余可读位数
} BitReader;
\end{lstlisting}
\end{itemize}

\subsection{主程序流程与模块层次}
\subsubsection{压缩程序 \texttt{compress.c}}
\begin{verbatim}
main
├── 获取输入/输出文件名（命令行或交互）
├── 统计频率 (fread 循环)
├── build_huffman_tree
│   ├── create_min_heap
│   ├── insert_min_heap
│   └── extract_min / heapify
├── generate_codes (递归遍历树)
├── 写入文件头 (fwrite)
├── 编码并写入比特流
│   ├── bit_writer_create
│   ├── bit_writer_write (使用 code_table)
│   └── bit_writer_flush
└── 释放所有资源
\end{verbatim}

\subsubsection{解压程序 \texttt{decompress.c}}
\begin{verbatim}
main
├── 获取输入文件/输出目录
├── 验证魔数，读取文件头
├── 读取频率表
├── rebuild_huffman_tree
│   ├── create_min_heap (变体)
│   └── extract_min / heapify (变体)
├── 构造输出路径并打开文件
├── 解码比特流
│   ├── bit_reader_create
│   ├── bit_reader_read (逐位返回)
│   └── 树遍历与字节输出
└── 释放所有资源
\end{verbatim}

两个程序各自独立编译，通过相同的频率表约定保证哈夫曼树的重建一致性。

\section{详细设计}

\subsection{核心算法伪代码}

\subsubsection{哈夫曼树构建（压缩与解压共用思路）}
\begin{lstlisting}
function build_huffman_tree(freq[256]) -> Node*:
    heap = create_min_heap(256)
    for symbol = 0 to 255:
        if freq[symbol] > 0:
            node = new Node(symbol, freq[symbol])
            insert_min_heap(heap, node)
    if heap.size == 0: return NULL
    while heap.size > 1:
        left = extract_min(heap)
        right = extract_min(heap)
        parent = new Node(-1, left.freq + right.freq)
        parent.left = left; parent.right = right
        insert_min_heap(heap, parent)
    root = extract_min(heap)
    销毁 heap->array 和 heap
    return root
\end{lstlisting}

\subsubsection{编码表生成（递归）}
\begin{lstlisting}
function generate_codes(root, arr[], depth, code_table[]):
    if root.left != NULL:
        arr[depth] = 0
        generate_codes(root.left, arr, depth+1, code_table)
    if root.right != NULL:
        arr[depth] = 1
        generate_codes(root.right, arr, depth+1, code_table)
    if root.left == NULL and root.right == NULL:  // 叶子
        code_table[root.symbol].len = depth
        将 arr[0..depth-1] 按位复制到 code_table[root.symbol].code 中
\end{lstlisting}

\subsubsection{位写入}
\begin{lstlisting}
function bit_writer_write(bw, code_table[], symbol):
    code = code_table[symbol].code
    len = code_table[symbol].len
    for i = 0 to len-1:
        bit = (code[i/8] >> (7 - i%8)) & 1
        if bit == 1: bw.buf |= (1 << (7 - bw.pos))
        bw.pos++
        if bw.pos == 8:
            fwrite(&bw.buf, 1, 1, bw.fp)
            bw.buf = 0; bw.pos = 0
    end
\end{lstlisting}

\subsubsection{解压主循环}
\begin{lstlisting}
while bytes_written < original_size:
    bit = bit_reader_read(br)   // 返回0/1或-1
    if bit < 0: error
    current = (bit == 0) ? current->left : current->right
    if current->left == NULL and current->right == NULL:
        fputc(current->symbol, fout)
        bytes_written++
        current = root
\end{lstlisting}

\subsection{函数调用关系图（压缩）}
\begin{verbatim}
main
├── build_huffman_tree
│   ├── create_min_heap / free_min_heap
│   ├── insert_min_heap
│   │   └── swap_nodes
│   ├── extract_min
│   │   └── heapify
│   └── free_min_heap
├── generate_codes
├── 文件写入操作 (fwrite / rewind)
├── bit_writer_create
├── bit_writer_write
│   └── fwrite
├── bit_writer_flush
└── free_tree
\end{verbatim}

\section{调试分析报告}

\subsection{调试过程中遇到的问题及解决方法}
\begin{enumerate}
    \item \textbf{问题1：Permission denied 错误} \\
    运行时输出 \texttt{无法创建压缩文件 xxx: Permission denied}，原因多为输出路径指向了只读目录或系统保护目录。\\
    \textbf{解决方法}：增加详细的错误输出（使用 \texttt{strerror(errno)}），并在程序启动时允许交互式输入文件名，默认在当前目录生成输出文件，避免权限问题。

    \item \textbf{问题2：空文件处理} \\
    若对空文件直接压缩，频率表全零，建树函数返回 \texttt{NULL}，后续编码读写会崩溃。\\
    \textbf{解决方法}：在统计完总字节数后检查 \texttt{total\_bytes == 0}，立即输出提示并退出，不进入压缩流程。

    \item \textbf{问题3：单一字节文件编码长度为 0} \\
    当文件只有一种字节时，该字节的哈夫曼编码可能被错误地设为 0 位，导致压缩文件不含有效编码，解压失败。\\
    \textbf{解决方法}：在建树逻辑中，若堆中仅一个节点，直接将其作为根节点返回；此时 \texttt{generate\_codes} 对该节点生成的编码长度为 1（根据需要设定），保证编码有效。经测试，现有实现已正确处理此种情况。
\end{enumerate}

\subsection{算法时空复杂度分析}
\begin{itemize}
    \item \textbf{压缩阶段}：
        \begin{itemize}
            \item 统计频率：一次遍历源文件，时间复杂度 \(O(N)\)，\(N\) 为文件字节数。
            \item 构建哈夫曼树：符号最多 256 个，堆操作复杂度 \(O(256 \log 256)\)，常数量级。
            \item 编码生成：递归遍历全树，同样常数时间。
            \item 编码写入：每个字节根据编码长度进行位写入，整体复杂度 \(O(N)\)。
            \item 空间复杂度：频率表、编码表、树节点等均占用常量级内存，位缓冲仅 1 字节，故空间开销为 \(O(1)\) 额外内存。
        \end{itemize}
    \item \textbf{解压阶段}：
        \begin{itemize}
            \item 读取头信息：常数时间。
            \item 重建哈夫曼树：与压缩相同，常数时间。
            \item 解码比特流：每处理一位沿树走一步，总位移次数等于压缩后的比特数，最坏情况下等于 \(O(N)\)（若编码长度远小于 8 位，可能略多，但总量级仍为 \(O(N)\)）。
            \item 空间复杂度：同样为 \(O(1)\) 额外内存（除原始文件数据外）。
        \end{itemize}
\end{itemize}

\subsection{改进设想}
\begin{itemize}
    \item 可以将多个文件打包成一个 \texttt{.huff} 存档。
    \item 在压缩、解压时增加进度显示，提升用户体验。
\end{itemize}

\section{用户使用说明}
直接运行对应可执行文件，按提示输入文件名：
\begin{enumerate}
    \item \textbf{压缩}：
        \begin{enumerate}
            \item 双击 \texttt{compress.exe} 或在命令行运行。
            \item 输入要压缩的源文件名，例如 \texttt{photo.jpg}。
            \item 输入压缩后的文件名，例如 \texttt{photo.huff}（可仅指定文件名，默认生成在当前目录）。
            \item 程序将显示 ``压缩完成：photo.jpg -> photo.huff''。
        \end{enumerate}
    \item \textbf{解压}：
        \begin{enumerate}
            \item 双击 \texttt{decompress.exe} 或在命令行运行。
            \item 输入要解压的压缩文件名，例如 \texttt{photo.huff}。
            \item 输入解压输出目录（直接回车则解压到当前目录）。
            \item 程序将显示 ``解压完成：photo.huff -> photo.jpg''。
        \end{enumerate}
\end{enumerate}

例如：
\begin{verbatim}
compress.exe test.png test.huff
decompress.exe test.huff .\output
\end{verbatim}

\section{测试结果}

\subsection{测试环境}
\begin{itemize}
    \item 操作系统：Windows 11
    \item 编译器：MinGW-W64 GCC 8.1.0
    \item 编译选项：\texttt{-std=c11 -Wall}（可选关警告）
\end{itemize}

\subsection{测试用例1：标准文本文件压缩解压}
\textbf{测试步骤}：
\begin{enumerate}
    \item 创建文本文件 \texttt{test.txt}，内容为 ``hello world this is a test file for huffman coding''。
    \item 执行 \texttt{compress.exe test.txt out.huff}。
    \item 执行 \texttt{decompress.exe out.huff}，得到还原文件 \texttt{test.txt}。
    \item 对比原文件和还原文件的 MD5 值。
\end{enumerate}
\textbf{结果}：压缩成功，解压后的文件与原文件完全一致，MD5 码相同。压缩文件大小比原文件大幅减小（本例中约减少 30\%）。

\subsection{测试用例2：二进制图片文件}
\textbf{测试数据}：一张 128×128 的 PNG 图片，大小 12,345 字节。\\
\textbf{操作}：压缩为 \texttt{img.huff}，再解压。\\
\textbf{结果}：复原图片与原图片逐字节相同，可正常打开查看。

\subsection{测试用例3：单一字节文件}
\textbf{测试数据}：文件 \texttt{single.bin} 包含 10,000 个字节 0x41。\\
\textbf{操作}：压缩后大小为几百字节（主要存储头信息和频率表），解压得到完全一致的 10,000 个 0x41。\\
\textbf{结果}：极限情况处理正确。

\subsection{测试用例4：空文件}
\textbf{测试数据}：一个 0 字节的空文件 \texttt{empty.dat}。\\
\textbf{操作}：运行压缩程序，程序立即输出“输入文件为空，无法压缩。”并返回错误码。\\
\textbf{结果}：正确处理边界情况。

\subsection{测试用例5：损坏的压缩文件}
\textbf{测试数据}：创建一个包含非 ``HUFF`` 魔数的文件。\\
\textbf{操作}：运行解压程序，输出“文件格式错误，不是有效的 Huffman 压缩文件。”\\
\textbf{结果}：错误检测有效。

\subsection{测试用例6：权限不足的输出路径}
\textbf{测试数据}：试图将压缩文件写入 \texttt{C:\textbackslash Windows\textbackslash System32\textbackslash test.huff}。\\
\textbf{操作}：运行压缩程序，系统返回“无法创建压缩文件 ... Permission denied”。\\
\textbf{结果}：出现错误提示，程序正常退出。

\end{document}